\clearpage\SourceLeaf{206}
\selectlanguage{german}\fancyhead[L]{\small ANZEIGE.}\pdfbookmark[0]{ANZEIGEN.}{anzeigen206}
\begin{center}\textbf{A N Z E I G E N.}\end{center}
\par\Needspace{9\baselineskip}\begin{center}
\rule{0.52\linewidth}{0.3pt}\par\medskip
Göttingische gelehrte Anzeigen.\quad 1812 Februar 10.\par\medskip
\rule{0.52\linewidth}{0.3pt}
\end{center}

Die logarithmischen und Kreisfunctionen, als die einfachsten Arten der
transscendenten Functionen, sind diejenigen, womit sich die Analysten am
meisten beschäftigt haben. Sie verdienten diese Ehre sowohl wegen ihres steten
Eingreifens in fast alle mathematische Untersuchungen, theoretische und practische,
als wegen des fast unerschöpflichen Reichthums an interessanten Wahrheiten, den
ihre Theorie darbietet\quad Weit weniger sind bisher andere transscendente Functionen
bearbeitet, die sich auf jene nicht zurückführen lassen, sondern als eigne
höhere Gattungen betrachtet werden müssen, die nur in speciellen Fällen mit
jenen zusammenhängen. Und doch sind manche solcher Functionen nicht minder
fruchtbar an interessanten Relationen, und daher dem, welcher die Analyse um
ihrer selbst willen ehrt, nicht minder wichtig; so wie ihr häufiges Vorkommen
bei mancherlei andern Untersuchungen sie dem empfehlen muss, der gern erst
nach practischem Nutzen fragt. Professor \textsc{Gauss} hat sich mit Untersuchungen
über dergleichen höhere transscendente Functionen schon seit vielen Jahren
beschäftigt. deren weit ausgedehnte Resultate das Gesagte bestätigen. Einen,
verhältnissmässig freilich nur sehr kleinen, Theil derselben, der gleichsam als
Einleitung zu einer künftig zu liefernden Reihe von Abhandlungen angesehen werden
kann, hat er am 30. Januar unter der Aufschrift

\begin{center}\emph{Disquisitiones generales circa seriem infinitam}\end{center}
\[
\begin{aligned}&1+\frac{\alpha.\beta}{1.\gamma}x
+\frac{\alpha.\alpha+1.\beta.\beta+1}{1.2.\gamma.\gamma+1}xx
\\[5pt]&+\frac{\alpha.\alpha+1.\alpha+2.\beta.\beta+1.\beta+2}{1.2.3.\gamma.\gamma+1.\gamma+2}x^3+\mathit{etc.}\end{aligned}
\]
\begin{center}\emph{Pars prior,}\end{center}
\SourceLeaf{207}
als eine Vorlesung der königl. Gesellschaft der Wissenschaften übergeben, woraus
wir hier die Hauptmomente des Inhalts anzeigen wollen.

Die transscendenten Functionen haben ihre wahre Quelle allemal. offen liegend
oder versteckt, im Unendlichen. Die Operationen des Integrirens, der Summationen
unendlicher Reihen, der Entwickelung unendlicher Producte, ins Unendliche
fortlaufender continuirlicher Brüche, oder überhaupt die Annäherung an
eine Grenze durch Operationen, die nach bestimmten Gesetzen ohne Ende fortgesetzt
werden --- diess ist der eigentliche Boden, auf welchem die transscendenten
Functionen erzeugt werden, oder wenn man lieber sich eines andern Bildes
bedienen will, diess sind die eigentlichen Wege, auf welchen man dazu gelangt.
Zu \emph{einem} Ziele führen gewöhnlich mehrere solcher Wege: die Umstände und die
Zwecke, welche man sich vorsetzt, müssen bestimmen, welchen man zuerst oder
vorzugsweise wählen will. Die Reihe, welche den Gegenstand gegenwärtiger Abhandlung
ausmacht, ist von einer sehr umfassenden Allgemeinheit. Sie stellt,
je nachdem die Grössen $\alpha$, $\beta$, $\gamma$ (welche nebst $x$ der Verfasser durch die Benennungen
erstes, zweites, drittes, viertes \emph{Element} unterscheidet, so wie er, Kürze
halber, die ganze durch die Reihe dargestellte Function mit den Zeichen $F(\alpha,\beta,\gamma,x)$
bezeichnet) so oder anders bestimmt werden, algebraische, logarithmische, trigonometrische
oder höhere transscendente Functionen dar, und man kann behaupten,
dass bisher kaum irgend eine transscendente Function von den Analysten untersucht
sei, die sich nicht auf diese Reihe zurückführen liesse. Eine grosse
Menge von Wahrheiten, welche in Beziehung auf solche schon in Betrachtung gezogene
Functionen schon aufgefunden sind, lassen sich aus der allgemeinen Natur
der durch unsere Reihen dargestellten Function ableiten, und schon um desswillen
würden Untersuchungen darüber die Aufmerksamkeit der Geometer verdienen,
obwohl diess nur als Nebensache, und die Eröffnung des Zuganges zu neuen
Wahrheiten als Hauptzweck zu betrachten ist. Gegenwärtige Abhandlung enthält
nur erst die Hälfte von den allgemeinen Untersuchungen des Verfassers, deren
zu grosser Umfang eine Theilung nothwendig machte. Hier gilt eben die
Reihe selbst als Ursprung der transscendenten Functionen, welche in dem weitern
Verfolg der Arbeit aus einer allgemeiner anwendbaren Quelle werden abgeleitet,
und aus einem höheren Gesichtspunkte betrachtet werden. Die erstere
Erzeugung macht, ihrer Natur nach, die Einschränkung auf die Fälle nothwendig,
wo die Reihe convergirt, also wo das vierte Element $x$. positiv oder negativ,
\SourceLeaf{208}
den Werth $1$ nicht überschreitet: die andere Erzeugungsart wird diese Beschränkung wegräumen. Allein eben die erstere Erzeugungsart führt schon zu einer Menge merkwürdiger Wahrheiten auf einem bequemern und gleichsam mehr elementarischen Wege, und desswegen hat der Verf. damit den Anfang gemacht.

Diese erste Hälfte der Untersuchungen zerfällt in drei Abschnitte, welchen einige allgemeine Bemerkungen voraus geschickt sind. Um eine Probe von der ausgedehnten Anwendbarkeit der Reihe zu geben, sind auf dieselbe zuvörderst $23$ verschiedene Reihenentwickelungen algebraischer, logarithmischer und trigonometrischer Functionen zurück geführt, so wie, als ein Beispiel höherer transcendenter Functionen, die Coëfficienten der aus der Entwickelung von $(aa+bb-2ab\cos\varphi)^n$ entspringenden, nach den Cosinus der Vielfachen von $\varphi$ fortschreitenden, Reihe, und zwar letztere auf drei verschiedene Arten.

Der \textit{erste Abschnitt} beschäftigt sich mit den Relationen zwischen solchen Reihen von der obigen Form, in welchen die Werthe eines der drei ersten Elemente um eine Einheit verschieden, die Werthe der drei übrigen hingegen gleich sind. Dergleichen Reihen nennt der Verf. \textit{series contiguae}, im Deutschen könnte man sie etwa \textit{verwandte Reihen} nennen. Jeder Reihe $F(\alpha,\beta,\gamma,x)$ stehen also sechs verwandte zur Seite, nämlich $F(\alpha+1,\beta,\gamma,x)$, $F(\alpha-1,\beta,\gamma,x)$, $F(\alpha,\beta+1,\gamma,x)$, $F(\alpha,\beta-1,\gamma,x)$, $F(\alpha,\beta,\gamma+1,x)$, $F(\alpha,\beta,\gamma-1,x)$, und es wird hier gezeigt, dass es zwischen der ersten und je zweien der verwandten eine lineare Gleichung gibt. Funfzehn Gleichungen entspringen auf diese Weise. Es folgt hieraus das wichtige Theorem, dass, wenn $\alpha'-\alpha$, $\alpha''-\alpha$, $\beta'-\beta$, $\beta''-\beta$, $\gamma'-\gamma$, $\gamma''-\gamma$, ganze Zahlen sind, auch zwischen $F(\alpha,\beta,\gamma,x)$, $F(\alpha',\beta',\gamma',x)$, $F(\alpha'',\beta'',\gamma'',x)$, eine lineare Gleichung Statt findet, und also, allgemein zu reden, aus den Werthen zweier dieser Functionen der Werth der dritten abgeleitet werden kann. Einige der einfachsten oder sonst merkwürdigen Fälle hat der Verf. hier noch besonders zusammengestellt.

Der \textit{zweite Abschnitt} gibt Verwandlungen in continuirliche Brüche, und zwar für die Quotienten
\[
\frac{F(\alpha,\beta+1,\gamma+1,x)}{F(\alpha,\beta,\gamma,x)},\quad
\frac{F(\alpha,\beta+1,\gamma,x)}{F(\alpha,\beta,\gamma,x)},\quad
\frac{F(\alpha-1,\beta+1,\gamma,x)}{F(\alpha,\beta,\gamma,x)}
\]
auf welche sich noch drei andere durch die offenbar verstattete Vertauschung der beiden ersten Elemente zurück führen lassen. Von diesen Lehrsätzen sind fast alle bisher bekannte Entwickelungen in continuirliche Brüche nur specielle Fälle.
\SourceLeaf{209}
Vorzüglich merkwürdig ist der Fall, wo man in der zweiten Entwickelung $\beta=0$ setzt. Es folgt daraus ein Lehrsatz, welchen wir seiner umfassenden Anwendbarkeit wegen hier beifügen. Die Function $F(\alpha,1,\gamma.x)$ oder, was einerlei ist, die Reihe
\[
1+\frac{\alpha}{\gamma}x
+\frac{\alpha\,.\,\alpha+1}{\gamma\,.\,\gamma+1}xx
+\frac{\alpha\,.\,\alpha+1\,.\,\alpha+2}{\gamma\,.\,\gamma+1\,.\,\gamma+2}x^3+\mathrm{etc.}
\]
gibt den continuirlichen Bruch
\[
\dfrac{1}{\displaystyle 1-\dfrac{ax}{\displaystyle 1-\dfrac{bx}{\displaystyle 1-\dfrac{cx}{\displaystyle 1-\dfrac{dx}{1-\mathrm{etc.}}}}}}
\]
wo die Coëfficienten $a,b,c,d$ etc. nach folgendem Gesetze fortschreiten:
\[
\begin{array}{ll}
a=\dfrac{\alpha}{\gamma}&b=\dfrac{\gamma-\alpha}{\gamma(\gamma+1)}\\[6pt]
c=\dfrac{(\alpha+1)\gamma}{(\gamma+1)(\gamma+2)}&d=\dfrac{2(\gamma+1-\alpha)}{(\gamma+2)(\gamma+3)}\\[6pt]
e=\dfrac{(\alpha+2)(\gamma+1)}{(\gamma+3)(\gamma+4)}&f=\dfrac{3(\gamma+2-\alpha)}{(\gamma+4)(\gamma+5)}\quad\mbox{u.\,s.\,f.}
\end{array}
\]
Hiernach lassen sich z.\,B. die Potenz eines Binomium, die Reihen für $\log(1+x)$, $\log\dfrac{1+x}{1-x}$, für Exponentialgrössen, für den Bogen durch die Tangente oder durch den Sinus u.\,a. in unendliche continuirliche Brüche verwandeln. Auch beruhen hierauf die in der \textit{Theoria motus corporum coelestium} gegebenen Verwandlungen in solche Brüche, deren Beweise hier von dem Verfasser nachgeholt werden.

Bei weitem den grössten Theil der Abhandlung nimmt der \textit{dritte Abschnitt} ein, in welchem von dem Werthe der Reihe gehandelt wird, wenn man das vierte Element $=1$ setzt. Nachdem zuvörderst mit geometrischer Schärfe bewiesen, dass die Reihe für $x=1$ nur dann zu einer endlichen Summe convergire, wenn $\gamma-\alpha-\beta$ eine positive Grösse ist, führt der Verf. diese Summe, oder $F(\alpha,\beta,\gamma,1)$ auf den Ausdruck $\dfrac{\Pi(\gamma-1)\,.\,\Pi(\gamma-\alpha-\beta-1)}{\Pi(\gamma-\alpha-1)\,.\,\Pi(\gamma-\beta-1)}$ zurück, wo die Charakteristik $\Pi$ eine eigene Art transcendenter Functionen andeutet, deren Erzeugung der Verf. auf ein unendliches Product gründet. Diese in der ganzen Analyse höchst wichtige Function ist im Grunde nichts anders als \textsc{Eulers} inexplicable Function
\[
\Pi z=1\,.\,2\,.\,3\,.\,4\,.\,.\,.\,.\,.\,z
\]
\SourceLeaf{210}
allein \textit{diese} Erzeugungsart oder Definition ist, nach des Verf. Urtheil, durchaus unstatthaft, da sie nur für ganze positive Werthe von $z$ einen klaren Sinn hat. Die vom Verf. gewählte Begründungsart ist allgemein anwendbar, und gibt selbst bei imaginären Werthen von $z$ einen eben so klaren Sinn, wie bei reellen, und man läuft dabei durchaus keine Gefahr, auf solche Paradoxen und Widersprüche zu gerathen, wie ehedem Hr.\ \textsc{Kramp} bei seinen numerischen Facultäten, die sich, wie man leicht zeigen kann, auf obige Function zurückführen lassen, aber zur Aufnahme in die Analyse weniger geeignet scheinen, als diese, da jene von drei Grössen abhängig sind, diese nur von Einer abhängt, und doch als eben so allgemein betrachtet werden muss. Der Verf. wünscht dieser transscendenten Function $\Pi z$ in der Analyse das Bürgerrecht gegeben zu sehen, wozu vielleicht die Wahl eines eigenen Namens für dieselbe am beförderlichsten sein würde: das Recht dazu mag demjenigen vorbehalten bleiben, der die wichtigsten Entdeckungen in der Theorie dieser der Anstrengungen der Geometer sehr würdigen Function machen wird. Hier ist von dem Verf. bereits eine bedeutende Anzahl merkwürdiger, sie betreffender, Theoreme zusammen gestellt, wovon ein Theil als neu zu betrachten ist. Der Raum verstattet uns nicht, in das Detail derselben hier einzugehen: nur das eine heben wir davon aus, dass der Werth des Integrals $\int x^{\lambda-1}(1-x^\mu)^\nu\,\mathrm{d}x$ von $x=0$ bis $x=1$ leicht auf die Function $\Pi$ zurückgeführt werden kann, und dass alle die von \textsc{Euler} für dergleichen Integrale zum Theil mühsam gefundenen Relationen sich mit grösster Leichtigkeit aus den allgemeinen Eigenschaften jener Functionen ableiten lassen, so wie umgekehrt allemal $\Pi z$, wenn $z$ eine Rationalgrösse ist, sich durch einige solche bestimmte Integrale darstellen lässt.

Nicht weniger merkwürdig ist die aus der Differentiation von $\Pi z$ entspringende, gleichfalls transscendente, Function, oder vielmehr
\[
\frac{\mathrm{d}\log\Pi z}{\mathrm{d}z}=\frac{\mathrm{d}\Pi z}{\Pi z\,.\,\mathrm{d}z}
\]
welche der Verfasser mit $\Psi z$ bezeichnet hat, und die gleichfalls eine besondere \textit{Benennung} verdiente. Von den zahlreichen merkwürdigen Eigenschaften dieser Function, welche in der Abhandlung aufgestellt sind, führen wir hier nur die Eine an, dass allgemein $\Psi z-\Psi 0$, wenn $z$ eine rationale Grösse ist, auf Logarithmen und Kreisfunctionen zurückgeführt werden kann; $\Psi 0$ selbst aber ist die bekannte, von \textsc{Euler} und Andern untersuchte, Zahl $0{,}5772156649\,.\,.\,.\,.$ negativ genommen, welche der Verfasser hier, nach einer von ihm selbst geführten Rech\SourceLeaf{211}%
nung, auf 23 Decimalen mittheilt, wovon die letzten von \textsc{Mascheroni}'s Bestimmung etwas abweichen. Uebrigens hängen sowohl $\Pi z$, als $\Psi z$, mit mehreren merkwürdigen Integralen für bestimmte Werthe der veränderlichen Grösse zusammen.

Diesem dritten Abschnitte ist noch eine unter der Aufsicht des Professors \textsc{Gauss} von Hrn.\ \textsc{Nicolai} mit grösster Sorgfalt berechnete Tafel für $\log\Pi z$ und für $\Psi z$ beigefügt, worin das Argument $z$ durch alle einzelnen Hunderttheile von 0 bis 1 fortschreitet; aus der Theorie dieser Functionen ist klar, dass man auf diese Werthe von $z$ alle andere leicht zurück führen kann.

\par\Needspace{9\baselineskip}\begin{center}
\rule{0.6\linewidth}{0.4pt}\par\medskip
Göttingische gelehrte Anzeigen.\quad 1814 September 26.\par\medskip
\rule{0.6\linewidth}{0.4pt}
\end{center}\pdfbookmark[1]{Göttingische gelehrte Anzeigen. 1814 September 26.}{anzeige1814}

Der königl.\ Societät wurde am 16.\ September von dem Prof.\ \textsc{Gauss} eine Vorlesung eingereicht überschrieben:

\begin{center}
\textit{Methodus nova integralium valores per approximationem inveniendi.}
\end{center}

Unter den verschiedenen Methoden zur genäherten Bestimmung der Integrale, oder wie es in der Sprache der ältern Analysten hiess, zur genäherten Quadratur krummliniger Figuren, ist die \textsc{Newton}-\textsc{Cotes}ische, welche sich auf die Interpolationsmethode gründet, eine der brauchbarsten. \textsc{Newton} hatte eine Auflösung der Aufgabe gegeben, durch eine beliebige Anzahl gegebener Punkte eine parabolische Curve zu ziehen, deren immer leicht ausführbare Quadratur dann näherungsweise die Stelle der Quadratur der eigentlich vorgegebenen durch jene Punkte gehenden Curve vertreten kann, und zwar desto genauer, je mehr Punkte man in Anwendung bringt. \textsc{Newton} hatte es indessen bei dieser allgemeinen Andeutung bewenden lassen und nur gleichsam beispielsweise für den Fall von vier in gleichen Zwischenräumen liegenden Ordinaten $A$, $B$, $C$, $D$ den genäherten Flächenraum zwischen der ersten und letzten, wenn deren Entfernung $=R$ ist, durch $(\frac{1}{8}A+\frac{3}{8}B+\frac{3}{8}C+\frac{1}{8}D)R$ angeführt. \textsc{Cotes}, welcher für sich, und noch ehe \textsc{Newtons} Schrift \textit{Methodus differentialis} erschienen war, schon im Jahre 1707 ähnliche Untersuchungen angestellt hatte, wurde durch die zierliche Form, in welcher \textsc{Newton} das Endresultat in obigem Beispiele dargestellt hatte (\textit{pulcherrima et}
